\documentclass[11pt,a4paper]{article}

% --- 核心宏包与中文支持 ---
\usepackage[utf8]{inputenc}
\usepackage{xeCJK}
\usepackage{amsmath, amssymb, amsfonts}
\usepackage{listings}
\usepackage{xcolor}
\usepackage{geometry}
\usepackage{tcolorbox}
\usepackage{booktabs}
\usepackage{tabularx}
\usepackage{setspace}
\usepackage{enumitem}
\usepackage[hidelinks]{hyperref}
\usepackage{bookmark}
\usepackage{fancyhdr}

% --- PDF 书签与元数据 ---
\hypersetup{
    colorlinks=false,
    pdfborder={0 0 0},
    bookmarks=true,
    bookmarksnumbered=true,
    bookmarksopen=true,
    bookmarksopenlevel=2,
    pdfcreator={XeLaTeX},
    pdftitle={08 - Pandas Library I},
    pdfauthor={黄子扬},
}

% --- 页面美化设置 ---
\geometry{top=0.7in,bottom=0.8in,left=0.6in,right=0.6in}
\onehalfspacing
\tcbuselibrary{skins,breakable}
\setlength{\parindent}{0pt}
\setlength{\parskip}{0.6em}

% --- 代码样式配置 ---
\definecolor{mainblue}{RGB}{0,51,102}
\definecolor{examred}{RGB}{180,0,0}
\definecolor{codebg}{RGB}{248,248,248}
\definecolor{commentgreen}{RGB}{0,128,0}

\lstset{
    backgroundcolor=\color{codebg},
    basicstyle=\ttfamily\footnotesize,
    keywordstyle=\color{blue}\bfseries,
    commentstyle=\color{commentgreen},
    stringstyle=\color{orange},
    breaklines=true,
    showstringspaces=false,
    frame=leftline,
    xleftmargin=2em,
    numbers=left,
    numbersep=5pt,
    numberstyle=\tiny\color{gray},
    language=Python,
    morekeywords={None, True, False, pd, np, DataFrame, Series, loc, iloc, at, iat, merge, join, groupby, import, as, for, in, if}
}

% --- 自定义教学框 ---
\newtcolorbox{concept}[1]{colback=blue!5, colframe=mainblue, fonttitle=\bfseries, title={{【Concept / 核心概念】#1}}, arc=3pt, breakable}
\newtcolorbox{warning}[1]{colback=red!5, colframe=examred, fonttitle=\bfseries, title={{【Warning / 避坑指南】#1}}, arc=3pt, breakable}
\newtcolorbox{examplebox}[1]{colback=green!2, colframe=green!40!black, fonttitle=\bfseries, title={{【Example / 深度实操案例】#1}}, arc=2pt, breakable}

% --- 页眉页脚 ---
\pagestyle{fancy}
\fancyhf{}
\fancyhead[L]{\small\bfseries 08 - Pandas Library I / Pandas 库 I}
\fancyhead[R]{\small\thepage}
\renewcommand{\headrulewidth}{0.5pt}

% --- 文档信息 ---
\title{\Huge \bfseries 08 - Pandas Library I}
\author{\Large \textbf{哈尔滨工业大学 \quad 黄子扬}}
\date{2026年6月8日}

\begin{document}

\maketitle
\tableofcontents
\newpage

% ======================================================================
\section{Introduction to Pandas / Pandas 简介}
% ======================================================================

\begin{concept}{What is Pandas? / 什么是 Pandas？}
    Pandas is one of the most important libraries of Python for data analysis. It is well suited for:
    \begin{itemize}[itemsep=0.3em]
        \item Tabular data with heterogeneously-typed columns (like SQL tables or Excel spreadsheets).
        \item Ordered and unordered time series data.
        \item Arbitrary matrix data with row and column labels.
        \item Any other form of observational / statistical data sets.
    \end{itemize}
    Install: \texttt{pip install pandas}. Import: \texttt{import pandas as pd}.
\end{concept}

% ======================================================================
\section{The Core Structures / Pandas 两大核心结构}
% ======================================================================

\begin{concept}{Series vs DataFrame / 一维 Series vs 二维 DataFrame}
    \begin{itemize}[itemsep=0.5em]
        \item \textbf{Series}：The 1-dimensional analog of the DataFrame. Think of it as a single column with labels. It is size-immutable — operations that change its size are not allowed. \\
        \textbf{Series} 是一维结构，类似带标签的单列数据。每个 DataFrame 的列本质上都是一个 Series。
        \item \textbf{DataFrame}：A tabular representation of data (rows and columns), similar to a spreadsheet or SQL table. It is \textbf{size-mutable} — data can be added or deleted. \\
        \textbf{DataFrame} 是二维表格结构，其大小可动态调整（增删行列）。
    \end{itemize}
\end{concept}

\begin{concept}{The Peculiarity of Index / 索引的特殊性}
    In Pandas, \textbf{Index} are about \textbf{labels}, not just performance (like in SQL). They help with:
    \begin{enumerate}[itemsep=0.3em]
        \item \textbf{Selection}：Label-based data access.
        \item \textbf{Automatic alignment}：When performing operations between two DataFrames or Series, Pandas automatically aligns data by index labels.
    \end{enumerate}
    在 Pandas 中，索引是关于"标签"的。它帮助在两个表格运算时自动对齐数据。

    Special Index types: \texttt{MultiIndex} (hierarchical), \texttt{DatetimeIndex} (datetimes), \texttt{Float64Index} (floats), \texttt{CategoricalIndex} (categories).
\end{concept}

% ======================================================================
\section{Creating DataFrames / 创建表格的多种方式}
% ======================================================================

\subsection{The DataFrame Constructor / DataFrame 构造函数}

\begin{concept}{Constructor Signature / 构造函数签名}
    \texttt{pd.DataFrame(data=None, index=None, columns=None, dtype=None, copy=False)}
    \begin{itemize}[itemsep=0.3em]
        \item \texttt{data}：Input — dict, list, set, ndarray, iterable, or DataFrame.
        \item \texttt{index}：(Optional) Row index list. Default: range of integers 0, 1, …, n.
        \item \texttt{columns}：(Optional) Column labels list. Default: range of integers 0, 1, …, n.
        \item \texttt{dtype}：(Optional) Data type for the whole DataFrame. By default, inferred from data.
    \end{itemize}
\end{concept}

\subsection{From Dict / 从字典创建}

\begin{examplebox}{Creating DataFrame from Dict / 字典 → DataFrame}
    Dict keys become column labels; dict values become column data.
\begin{lstlisting}
import pandas as pd

# Simple dict / 简单字典
student_dict = {'Name': ['Joe', 'Nat'], 'Age': [20, 21], 'Marks': [85.10, 77.80]}
student_df = pd.DataFrame(student_dict)
print(student_df)
#    Name  Age  Marks
# 0   Joe   20  85.10
# 1   Nat   21  77.80

# With custom indices / 自定义索引
df = pd.DataFrame({'A': [1, 2, 3], 'B': [True, True, False],
                   'C': [0.496714, -0.138264, 0.647689]},
                  index=['i', 'ii', 'iii'])
print(df)
\end{lstlisting}
\end{examplebox}

\subsection{From List / 从列表创建}

\begin{examplebox}{Creating DataFrame from List / 列表 → DataFrame}
    By default, all list elements are added as a \textbf{single row}.
\begin{lstlisting}
# Simple list — becomes one row / 列表变单行
fruits_list = ['Apple', 10, 'Orange', 55.50]
fruits_df = pd.DataFrame(fruits_list)
print(fruits_df)

# With custom column names and index / 自定义列名和索引
fruits_list = ['Apple', 'Banana', 'Orange', 'Mango']
fruits_df = pd.DataFrame(fruits_list, columns=['Fruits'],
                         index=['Fruit1', 'Fruit2', 'Fruit3', 'Fruit4'])
print(fruits_df)

# Multi-dimensional list / 多维列表
fruits_list = [['Apple', 'Banana', 'Orange', 'Mango'], [120, 40, 80, 500]]
fruits_df = pd.DataFrame(fruits_list)
print(fruits_df)        # Each sublist becomes a row

# Transpose to swap rows/columns / 转置行列互换
fruits_df = pd.DataFrame(fruits_list).transpose()
print(fruits_df)
\end{lstlisting}
\end{examplebox}

\subsection{From CSV / 从 CSV 文件创建}

\begin{examplebox}{Reading CSV Files / 读取 CSV 文件}
    \texttt{pd.read\_csv("file.csv")} is the most common way to load large datasets.
\begin{lstlisting}
data1 = pd.read_csv("data/sample.csv")
print(type(data1))              # <class 'pandas.core.frame.DataFrame'>
print(data1.iloc[:, 2])         # Select 3rd column by position
\end{lstlisting}
\end{examplebox}

% ======================================================================
\section{Indexing \& Slicing / 索引与切片（期末核心考点）}
% ======================================================================

\begin{warning}{Label-based (.loc) vs Position-based (.iloc)}
    这是 Pandas 初学者最容易混淆的地方，务必背诵以下区别：
    \begin{itemize}[itemsep=0.8em]
        \item \textbf{\texttt{.loc[row\_labels, column\_labels]}}：Based on \textbf{Labels}. Slicing is \textbf{INCLUSIVE} of the right boundary! \\
        例：\texttt{df.loc['i':'iii']} 返回包含 \texttt{'iii'} 在内的所有行。
        \item \textbf{\texttt{.iloc[row\_positions, column\_positions]}}：Based on \textbf{Integer Positions}. Slicing is \textbf{EXCLUSIVE} of the right boundary (standard Python behavior). \\
        例：\texttt{df.iloc[0:2]} 只返回第 0 和第 1 行，\textbf{不含}第 2 行。
    \end{itemize}
\end{warning}

\begin{examplebox}{Indexing Examples / 索引实战}
\begin{lstlisting}
df = pd.DataFrame({'A': [1, 2, 3], 'B': [True, True, False],
                   'C': [0.496714, -0.138264, 0.647689]},
                  index=['i', 'ii', 'iii'])

# Column Selection / 列选择
print(df['A'])                      # Single column → returns Series
print(df[['A', 'C']])               # Multiple columns → returns DataFrame

# Row Selection with .loc (label-based)
print(df.loc[['i', 'iii']])         # Select rows 'i' and 'iii'
print(df.loc['i':'iii'])            # INCLUSIVE of 'iii'!

# Row Selection with .iloc (position-based)
print(df.iloc[[0, 1]])              # Select rows 0 and 1
print(df.iloc[:2])                  # EXCLUSIVE of index 2

# Combined row + column selection / 行列同时选择
print(df.loc['ii', 'C'])            # Scalar value at row 'ii', column 'C'
print(df.loc['i':'ii', ['A', 'C']]) # Rows 'i' to 'ii', columns A and C

# Scalar access / 单值快速访问
print(df.at[0, 'Name'])             # Fast label-based single value
print(df.iat[0, 0])                 # Fast position-based single value
\end{lstlisting}
\end{examplebox}

\begin{concept}{Column Selection Shortcuts / 列选择快捷方式}
    \begin{itemize}[itemsep=0.3em]
        \item \texttt{df['A']} — Single column, returns \textbf{Series}.
        \item \texttt{df[['A', 'C']]} — Multiple columns, returns \textbf{DataFrame}.
        \item \texttt{df.A} — Attribute lookup (works only if column name has no spaces and doesn't conflict with DataFrame methods).
    \end{itemize}
\end{concept}

% ======================================================================
\section{Series / Series 详解}
% ======================================================================

\begin{examplebox}{Creating and Selecting Series / Series 的创建与选择}
\begin{lstlisting}
# Three ways to select a column as Series / 三种选取列的方式
df['A']                 # __getitem__ style
df.loc[:, 'A']          # .loc style (all rows, column 'A')
df.A                    # Attribute lookup style

# Create a Series directly / 直接创建 Series
s = pd.Series([2, 4, 6, 8, 10])
print(s)
# 0    2
# 1    4
# 2    6
# 3    8
# 4    10
# dtype: int64
\end{lstlisting}
    Series share many of the same methods as DataFrames.
\end{examplebox}

% ======================================================================
\section{Inspecting Metadata / 查看元数据与统计}
% ======================================================================

\begin{concept}{Key Inspection Functions / 核心检查函数}
    \begin{itemize}[itemsep=0.8em]
        \item \textbf{\texttt{df.info()}}：Gives index range, column list, non-null counts, data types, and memory usage.
        \item \textbf{\texttt{df.describe()}}：Summary statistics (count, mean, std, min, 25\%, 50\%, 75\%, max) for \textbf{numerical columns only}.
        \item \textbf{\texttt{df.shape}}：Returns tuple \texttt{(rows, columns)}.
    \end{itemize}
\end{concept}

\begin{examplebox}{Metadata and Statistics / 元数据与统计示例}
\begin{lstlisting}
data1 = pd.read_csv("data/sample.csv")
data1.info()            # Full metadata summary
data1.describe()        # Statistics for numerical columns only
print(data1.shape)      # (rows, columns)
\end{lstlisting}
\end{examplebox}

\subsection{DataFrame Attributes / DataFrame 属性一览}

\begin{table}[h]
\centering
\renewcommand{\arraystretch}{1.3}
\begin{tabular}{@{}ll@{}}
\toprule
\textbf{Attribute / 属性} & \textbf{Description / 描述} \\ \midrule
\texttt{df.index}   & Range of the row index / 行索引范围 \\
\texttt{df.columns} & List of column labels / 列标签列表 \\
\texttt{df.dtypes}  & Column names and their data types / 列名及数据类型 \\
\texttt{df.values}  & All rows as numpy array / 所有数据转为 numpy 数组 \\
\texttt{df.empty}   & Check if DataFrame is empty / 检查是否为空 \\
\texttt{df.size}    & Total number of values (rows × columns) / 值的总数 \\
\texttt{df.shape}   & Number of rows and columns / 行列数元组 \\ \bottomrule
\end{tabular}
\end{table}

% ======================================================================
\section{DataFrame Selection / 数据选择函数}
% ======================================================================

\begin{examplebox}{head, tail, at, iat, get / 选择函数}
\begin{lstlisting}
student_dict = {'Name': ['Joe', 'Nat', 'Harry'], 'Age': [20, 21, 19],
                'Marks': [85.10, 77.80, 91.54]}
student_df = pd.DataFrame(student_dict)

# Top N rows / 前 N 行
student_df.head(2)

# Bottom N rows / 后 N 行
student_df.tail(2)

# Single value by label / 按标签获取单值
student_df.at[0, 'Name']            # 'Joe'

# Single value by position / 按位置获取单值
student_df.iat[0, 0]                # 'Joe'

# Get column / 获取列
student_df.get('Name')              # Returns Series
\end{lstlisting}
\end{examplebox}

% ======================================================================
\section{Data Manipulation / 数据修改操作}
% ======================================================================

\subsection{Inserting Columns / 插入列}

\begin{examplebox}{df.insert() / 插入列}
    \texttt{df.insert(loc=col\_position, column=new\_col\_name, value=default\_value)}
\begin{lstlisting}
student_dict = {'Name': ['Joe', 'Nat', 'Harry'], 'Age': [20, 21, 19],
                'Marks': [85.10, 77.80, 91.54]}
student_df = pd.DataFrame(student_dict)

# Insert new column 'Class' at position 2 with default value 'A'
student_df.insert(loc=2, column="Class", value='A')
print(student_df)
\end{lstlisting}
\end{examplebox}

\subsection{Dropping Columns / 删除列}

\begin{examplebox}{df.drop() / 删除列}
    \texttt{df.drop(columns=[col1, col2, ...])}
\begin{lstlisting}
# Delete the 'Age' column
student_df = student_df.drop(columns='Age')
print(student_df)
\end{lstlisting}
\end{examplebox}

\subsection{Conditional Update with .where() / 条件替换}

\begin{concept}{.where() Logic / .where() 替换逻辑}
    \texttt{df.where(condition, other=new\_value)}：
    \begin{itemize}[itemsep=0.3em]
        \item If \texttt{condition} is \textbf{True} → keep the original value.
        \item If \texttt{condition} is \textbf{False} → replace with \texttt{other}.
    \end{itemize}
    条件为假则替换，条件为真则保持不变。
\end{concept}

\begin{examplebox}{Conditional Replacement / 条件替换示例}
\begin{lstlisting}
student_dict = {'Name': ['Joe', 'Nat', 'Harry'], 'Age': [20, 21, 19],
                'Marks': [85.10, 77.80, 91.54]}
student_df = pd.DataFrame(student_dict)

# Replace marks < 80 with 0
filter = student_df['Marks'] > 80
student_df['Marks'].where(filter, other=0, inplace=True)
print(student_df)
# Joe: 85.10 stays (True), Nat: 77.80 → 0 (False), Harry: 91.54 stays (True)
\end{lstlisting}
\end{examplebox}

\subsection{Filtering Columns by Label / 按列标签过滤}

\begin{warning}{.filter() filters LABELS, NOT data! / .filter() 过滤的是标签而非数据}
    \texttt{df.filter(like='N', axis='columns')} applies the filter to the \textbf{column labels or row index}, NOT the actual data inside the DataFrame. \\
    \texttt{.filter()} 过滤的是"列名/行名"，而不是数据内容！
\end{warning}

\begin{examplebox}{.filter() Example / .filter() 示例}
\begin{lstlisting}
student_df = pd.DataFrame({'Name': ['Joe', 'Nat', 'Harry'],
                           'Age': [20, 21, 19],
                           'Marks': [85.10, 77.80, 91.54]})

# Only include columns whose label starts with 'N'
student_df = student_df.filter(like='N', axis='columns')
print(student_df)   # Only 'Name' column remains
\end{lstlisting}
\end{examplebox}

\subsection{Renaming Columns / 重命名列}

\begin{examplebox}{df.rename() / 列重命名}
    \texttt{df.rename(columns=\{'old\_name': 'new\_name'\})}
\begin{lstlisting}
student_df = student_df.rename(columns={'Marks': 'Percentage'})
print(student_df)
\end{lstlisting}
\end{examplebox}

% ======================================================================
\section{Combining Data / 数据合并 (Join \& Merge)}
% ======================================================================

\subsection{Join / 横向拼接}

\begin{examplebox}{df.join() / 横向拼接}
    \texttt{df1.join(df2)} joins DataFrames by index.
\begin{lstlisting}
student_df = pd.DataFrame({'Name': ['Joe', 'Nat'], 'Age': [20, 21]})
marks_df = pd.DataFrame({'Marks': [85.10, 77.80]})
joined_df = student_df.join(marks_df)
print(joined_df)
\end{lstlisting}
\end{examplebox}

\subsection{Merge / SQL 风格合并}

\begin{concept}{SQL-style Merging / SQL 风格合并}
    \texttt{pd.merge(df1, df2, how='...', on='column\_name')}:
    \begin{itemize}[itemsep=0.5em]
        \item \textbf{\texttt{how='inner'}}：Only keep rows where the key exists in \textbf{BOTH} tables. / 只保留两表都有的键。
        \item \textbf{\texttt{how='left'}}：Keep all rows from \texttt{df1}, fill missing \texttt{df2} data with \texttt{NaN}. / 保留左表所有行。
        \item \textbf{\texttt{how='right'}}：Keep all rows from \texttt{df2}, fill missing \texttt{df1} data with \texttt{NaN}. / 保留右表所有行。
        \item \textbf{\texttt{how='outer'}}：Keep all rows from \textbf{both} tables. / 保留两表所有行。
    \end{itemize}
\end{concept}

\begin{examplebox}{Merge Examples / 合并示例}
\begin{lstlisting}
df1 = pd.DataFrame({'Name': ['Joe', 'Nat', 'Davis'], 'Age': [20, 21, 23]})
df2 = pd.DataFrame({'Name': ['Joe', 'Nat', 'Alice'], 'Marks': [85.10, 77.80, 90.00]})

print(pd.merge(df1, df2, how='inner', on='Name'))   # Joe, Nat only
print(pd.merge(df1, df2, how='left', on='Name'))    # Joe, Nat, Davis
print(pd.merge(df1, df2, how='right', on='Name'))   # Joe, Nat, Alice
print(pd.merge(df1, df2, how='outer', on='Name'))   # Joe, Nat, Davis, Alice
\end{lstlisting}
\end{examplebox}

% ======================================================================
\section{Aggregation \& Sorting / 分组聚合与排序}
% ======================================================================

\subsection{GroupBy / 分组聚合}

\begin{concept}{Split-Apply-Combine / 拆分-应用-组合}
    \texttt{df.groupby(col\_label).mean()} splits data by groups, applies the aggregation, and combines results.
\end{concept}

\begin{examplebox}{GroupBy Example / 分组聚合示例}
\begin{lstlisting}
student_df = pd.DataFrame({'Name': ['Joe', 'Nat', 'Harry'],
                           'Class': ['A', 'B', 'A'],
                           'Marks': [85.10, 77.80, 91.54]})

# Average marks per class / 每个班级的平均分
result = student_df.groupby('Class').mean()
print(result)
#         Marks
# Class
# A      88.32
# B      77.80

# Reset index to make 'Class' a column again
result = result.reset_index()
\end{lstlisting}
\end{examplebox}

\subsection{Sorting / 排序}

\begin{examplebox}{df.sort\_values() / 排序}
    \texttt{df.sort\_values(by=['col1', 'col2'], ascending=[True, False])}
\begin{lstlisting}
student_df = pd.DataFrame({'Name': ['Joe', 'Nat', 'Harry'],
                           'Age': [20, 21, 19],
                           'Marks': [85.10, 77.80, 91.54]})

# Sort by 'Marks' ascending
student_df = student_df.sort_values(by=['Marks'])
# Sort by 'Marks' descending
student_df = student_df.sort_values(by=['Marks'], ascending=False)
\end{lstlisting}
\end{examplebox}

\subsection{Iteration / 遍历}

\begin{examplebox}{df.iterrows() / 逐行遍历}
    Returns \texttt{(index, row)} pairs for each row.
\begin{lstlisting}
student_df = pd.DataFrame({'Name': ['Joe', 'Nat'], 'Age': [20, 21], 'Marks': [85, 77]})

for index, row in student_df.iterrows():
    print(index, row)
# 0 Name     Joe
#   Age       20
#   Marks     85
#   Name: 0, dtype: object
\end{lstlisting}
\end{examplebox}

% ======================================================================
\section{Conversion / 格式转换}

\begin{examplebox}{df.to\_dict() / 转回字典}
\begin{lstlisting}
dict_result = student_df.to_dict()
print(dict_result)
# {'Name': {0: 'Joe', 1: 'Nat'}, 'Age': {0: 20, 1: 21}, ...}
\end{lstlisting}
\end{examplebox}

% ======================================================================
\section{Quick Reference / 速查表}
% ======================================================================

\begin{table}[h]
\centering
\small
\renewcommand{\arraystretch}{1.3}
\begin{tabular}{@{}ll@{}}
\toprule
\textbf{Operation / 操作} & \textbf{Syntax / 语法} \\ \midrule
创建 DataFrame(字典) & \texttt{pd.DataFrame(dict)} \\
创建 DataFrame(列表) & \texttt{pd.DataFrame(list)} \\
读取 CSV & \texttt{pd.read\_csv('file.csv')} \\
选择列 & \texttt{df['col']} / \texttt{df[['c1','c2']]} \\
标签索引 & \texttt{df.loc[row, col]} \\
位置索引 & \texttt{df.iloc[row, col]} \\
单值(标签) & \texttt{df.at[row, col]} \\
单值(位置) & \texttt{df.iat[row, col]} \\
元数据 & \texttt{df.info()} \\
统计摘要 & \texttt{df.describe()} \\
形状 & \texttt{df.shape} \\
前 N 行 & \texttt{df.head(n)} \\
后 N 行 & \texttt{df.tail(n)} \\
插入列 & \texttt{df.insert(loc, name, value)} \\
删除列 & \texttt{df.drop(columns=[...])} \\
条件替换 & \texttt{df.where(cond, other=x)} \\
重命名 & \texttt{df.rename(columns=\{old:new\})} \\
合并(SQL) & \texttt{pd.merge(df1, df2, how='inner', on='key')} \\
横向拼接 & \texttt{df1.join(df2)} \\
分组聚合 & \texttt{df.groupby('col').mean()} \\
排序 & \texttt{df.sort\_values('col')} \\
遍历 & \texttt{df.iterrows()} \\
转字典 & \texttt{df.to\_dict()} \\ \bottomrule
\end{tabular}
\end{table}

\end{document}
